% Part: first-order-logic
% Chapter: syntax-and-semantics
% Section: subformulas

\documentclass[../../../include/open-logic-section]{subfiles}

\begin{document}

\olfileid{fol}{syn}{sbf}

\olsection{\printtoken{P}{subformula}}

\begin{explain}
આપેલા !!{formula}ને ``બનાવતાં'' !!{formula}s વિશે વાત કરવી ઘણી વાર
ઉપયોગી છે. આપણે તેમને તેનાં \emph{!!{subformula}s} કહીએ છીએ. દરેક
!!{formula} પોતાનું !!a{subformula} ગણાય છે; $!A$થી જુદું એવું $!A$નું
ઉપસૂત્ર \emph{ઉચિત !!{subformula}} કહેવાય છે.
\end{explain}

\begin{defn}[તત્કાલિક !!^{subformula}]
જો $!A$ !!a{formula} હોય, તો $!A$નાં \emph{તત્કાલિક
!!{subformula}s} નીચે પ્રમાણે અનુમાનાત્મક રીતે વ્યાખ્યાયિત થાય છે:
\begin{enumerate}
\item આણ્વિક !!{formula}sનાં કોઈ તત્કાલિક !!{subformula}s નથી.

\tagitem{prvNot}{\indcase{!A}{\lnot !B}{$\indfrm$નું એકમાત્ર
    તત્કાલિક !!{subformula}~$!B$ છે.}}{}

\item \indcase{!A}{(!B \ast !C)}{$\indfrm$નાં તત્કાલિક
  !!{subformula}s $!B$ અને $!C$ છે ($\ast$ કોઈ પણ દ્વિસ્થાની સંયોજક
  છે).}

\tagitem{prvAll}{\indcase{!A}{\lforall[x][!B]}{$\indfrm$નું એકમાત્ર
    તત્કાલિક !!{subformula}~$!B$ છે.}}{}

\tagitem{prvEx}{\indcase{!A}{\lexists[x][!B]}{$\indfrm$નું એકમાત્ર
    તત્કાલિક !!{subformula}~$!B$ છે.}}{}
\end{enumerate}
\end{defn}

\begin{defn}[ઉચિત !!^{subformula}]
જો $!A$ !!a{formula} હોય, તો $!A$નાં \emph{ઉચિત !!{subformula}s}
નીચે પ્રમાણે પુનરાવર્તી રીતે વ્યાખ્યાયિત થાય છે:
\begin{enumerate}
\item આણ્વિક !!{formula}sનાં કોઈ ઉચિત !!{subformula}s નથી.

\tagitem{prvNot}{\indcase{!A}{\lnot !B}{$\indfrm$નાં ઉચિત
    !!{subformula}sમાં~$!B$ અને $!B$નાં બધાં ઉચિત !!{subformula}s
    આવે છે.}}{}

\item \indcase{!A}{(!B \ast !C)}{$\indfrm$નાં ઉચિત !!{subformula}sમાં
  $!B$, $!C$ તથા $!B$ અને $!C$નાં બધાં ઉચિત !!{subformula}s આવે છે.}

\tagitem{prvAll}{\indcase{!A}{\lforall[x][!B]}{$\indfrm$નાં ઉચિત
    !!{subformula}sમાં~$!B$ અને $!B$નાં બધાં ઉચિત !!{subformula}s
    આવે છે.}}{}

\tagitem{prvEx}{\indcase{!A}{\lexists[x][!B]}{$\indfrm$નાં ઉચિત
    !!{subformula}sમાં~$!B$ અને $!B$નાં બધાં ઉચિત !!{subformula}s
    આવે છે.}}{}
\end{enumerate}
\end{defn}

\begin{defn}[!!^{subformula}]
$!A$નાં !!{subformula}s એટલે $!A$ પોતે અને તેનાં બધાં ઉચિત
!!{subformula}s.
\end{defn}

\begin{explain}
તત્કાલિક !!{subformula}s અને ઉચિત !!{subformula}s આપણે જે રીતે
વ્યાખ્યાયિત કર્યાં છે તેમાં રહેલો સૂક્ષ્મ તફાવત નોંધો. પહેલા કિસ્સામાં
દરેક શક્ય સ્વરૂપના સૂત્ર~$!A$ માટે આપણે $!A$નાં તત્કાલિક
!!{subformula}s સીધાં વ્યાખ્યાયિત કર્યાં છે. આ કિસ્સાઓ દ્વારા આપેલી સ્પષ્ટ વ્યાખ્યા
છે અને તેના કિસ્સા !!{formula}sના ગણની અનુમાનાત્મક વ્યાખ્યાને અનુસરે
છે. બીજા કિસ્સામાં પણ બધાં !!{formula}sનો ગણ જે રીતે વ્યાખ્યાયિત થાય
છે તેને અનુસર્યા છીએ, પણ દરેક કિસ્સામાં એ ઉચિત !!{subformula}sને પણ
સમાવ્યાં છે જે નાનાં !!{formula}s $!B$, $!C$નાં છે; આ !!{formula}s
પોતે તો સામેલ છે જ. તેનાથી
વ્યાખ્યા \emph{પુનરાવર્તી} બને છે. સામાન્ય રીતે, અનુમાનાત્મક રીતે
વ્યાખ્યાયિત ગણ પરના વિધેયની વ્યાખ્યા (અહીં !!{formula}s પરની)
ત્યારે પુનરાવર્તી કહેવાય જ્યારે વિધેયની વ્યાખ્યાના કિસ્સાઓમાં વિધેયનો
પોતાનો ઉપયોગ થતો હોય. જોકે વ્યાખ્યા સુવ્યાખ્યાયિત રહે એ માટે ખાતરી
કરવી પડે કે આપણે હંમેશાં માત્ર એ દલીલો માટેનાં વિધેયનાં મૂલ્યો વાપરીએ
છીએ જે વ્યાખ્યાયિત થતી દલીલની ``પહેલાં'' આવે છે---અહીં
$(!B\ast !C)$નું ``ઉચિત !!{subformula}'' વ્યાખ્યાયિત કરતાં આપણે માત્ર
ઉચિત !!{subformula}s વાપરીએ છીએ જે ``અગાઉનાં'' !!{formula}s $!B$ અને
$!C$નાં છે.
\end{explain}

\begin{prop}
\ollabel{prop:subfrm-trans}
ધારો કે $!B$, $!A$નું ઉપસૂત્ર છે અને $!C$, $!B$નું ઉપસૂત્ર છે. તો
$!C$, $!A$નું ઉપસૂત્ર છે. બીજા શબ્દોમાં, ઉપસૂત્ર-સંબંધ પરંપરિત છે.
\end{prop}

\begin{prob}
\olref[fol][syn][sbf]{prop:subfrm-trans} સાબિત કરો.
\end{prob}

\begin{prop}
\ollabel{prop:count-subfrms}
ધારો કે $!A$ એ $n$ સંયોજકો અને પરિમાણકો ધરાવતું સૂત્ર છે. તો $!A$ને
વધુમાં વધુ $2n+1$ ઉપસૂત્રો છે.
\end{prop}

\begin{prob}
\olref[fol][syn][sbf]{prop:count-subfrms} સાબિત કરો.
\end{prob}

\end{document}
